../cictikz/src/cictikz/data/tex/cictikz_preamble.tex

% What gets added to the square law below about 100 nm.
%
% Replaces a figure reproduced from Lewyn's paper, which showed every
% effect at once. This draws four families and only four, because the
% sections that follow take them one at a time and the job of this
% figure is to say what is coming.
%
% Colour follows the house rules: armygreen for the mechanical, blue for
% fields, red for the two that do damage.

\begin{tikzpicture}[scale=0.9]
  % Shared pieces for the lr0_mosfet cross-section cartoons.
%
% Two families use this file:
%  - the field-effect intro (mosfet_off/subthreshold/strong_inversion...):
%    white background, glowing blue n+ blocks, gate box, S/G/D terminals
%  - the carrier cartoons (accumulated/depleted/weakinv/vds_*): gray
%    substrate, n+ boxes with blue minus signs, red plus signs in the
%    p- bulk, teal oxide strip under the gate
%
% The colour language follows the original hand drawings: free electrons
% are blue, holes are red, the lecture prose ("imagine the free electrons
% as a blue fluorescent color") depends on it.

\definecolor{echarge}{RGB}{30,60,220}   % electron blue
\definecolor{hcharge}{RGB}{220,30,30}   % hole red
\definecolor{oxteal}{RGB}{60,140,160}   % oxide/channel strip

% ------------------------------------------------------------------
% Family 1: field-effect intro frame.
% Geometry: n+ blocks [0,2.9] and [5.9,8.8] x [-1.5,0], gate [2.6,6.2].
% ------------------------------------------------------------------
\newcommand{\fetblock}[2]{% x0 x1: glowing blue block below the surface
  \fill[echarge!30, rounded corners=8pt] (#1-0.12,-1.62) rectangle (#2+0.12,0.10);
  \fill[echarge, rounded corners=6pt] (#1,-1.5) rectangle (#2,0);
}

\newcommand{\fetterminal}[3]{% x y label: stem with a small open circle
  \draw[thick] (#1,#2) -- (#1,#2+0.35);
  \draw[thick] (#1,#2+0.44) circle (0.09);
  \node[anchor=south] at (#1,#2+0.62) {#3};
}

\newcommand{\fetframe}{%
  \fetblock{0}{2.9}
  \fetblock{5.9}{8.8}
  % gate: rounded box sitting just above the surface
  \draw[line width=1.6pt, rounded corners=10pt] (2.6,0.12) rectangle (6.2,1.05);
  \node[anchor=center] at (4.4,0.58) {GATE};
  \fetterminal{1.3}{0}{S}
  \fetterminal{4.4}{1.05}{G}
  \fetterminal{7.4}{0}{D}
}

% ------------------------------------------------------------------
% Family 2: carrier cartoon frame.
% Surface at y=0. Substrate down to y=-3.6, x in [-0.4,10.4].
% n+ boxes [0.3,2.9] and [7.1,9.7] x [-1.15,0]. Gate [3.0,7.0].
% ------------------------------------------------------------------
\newcommand{\eminus}[2]{\node[echarge, anchor=center, scale=0.7] at (#1,#2) {$\boldsymbol{-}$};}
\newcommand{\hplus}[2]{\node[hcharge, anchor=center, scale=0.8] at (#1,#2) {$\boldsymbol{+}$};}

% rows of minus signs filling an n+ box
\newcommand{\efill}[2]{% x0 x1
  \foreach \y in {-0.25,-0.55,-0.85}{
    \foreach \x in {0.15,0.55,...,2.4}{
      \eminus{#1+\x}{\y}
    }
  }
}

\newcommand{\bulkframe}{%
  % substrate outline
  \draw[line width=2.2pt, gray!60, rounded corners=4pt]
    (-0.4,0) -- (10.4,0) -- (10.4,-3.6) -- (-0.4,-3.6) -- cycle;
  % n+ wells
  \draw[line width=1.4pt, rounded corners=3pt] (0.3,0) -- (0.3,-1.15) -- (2.9,-1.15) -- (2.9,0);
  \draw[line width=1.4pt, rounded corners=3pt] (7.1,0) -- (7.1,-1.15) -- (9.7,-1.15) -- (9.7,0);
  \efill{0.3}{2.9}
  \efill{7.1}{9.7}
  \node[anchor=south west, scale=0.8] at (0.32,-1.12) {$n+$};
  \node[anchor=south east, scale=0.8] at (9.68,-1.12) {$n+$};
  % deep bulk holes
  \foreach \y/\xoff in {-3.3/0.0, -2.85/0.25}{
    \foreach \x in {0.1,0.7,...,10.2}{
      \hplus{\x+\xoff}{\y}
    }
  }
  \node[hcharge!70!black, anchor=center] at (5.0,-3.95) {$p-$};
  % oxide strip and gate box
  \fill[oxteal!80] (3.0,0.02) rectangle (7.0,0.24);
  \draw[line width=1.6pt, rounded corners=6pt] (2.75,0.26) rectangle (7.25,1.15);
}

% one row of holes between the wells
\newcommand{\neckrow}[2]{% y xoff
  \foreach \x in {3.2,3.8,...,6.9}{
    \hplus{\x+#2}{#1}
  }
}

% holes to the sides underneath the wells
\newcommand{\sideholes}{%
  \foreach \y/\xoff in {-1.6/0.1, -2.15/0.35}{
    \foreach \x in {-0.1,0.5,...,2.6}{ \hplus{\x+\xoff}{\y} }
    \foreach \x in {7.3,7.9,...,10.1}{ \hplus{\x+\xoff}{\y} }
  }
}

% terminals with ground symbols, and the p- label
\newcommand{\bulkterminalgnd}[3]{% x label side(-1 left / 1 right)
  \draw[thick] (#1,0) -- (#1,0.5) -- ({#1+#3*0.7},0.5);
  \draw[thick] ({#1+#3*0.7},0.35) -- ({#1+#3*0.7},0.65);
  \draw[thick] ({#1+#3*0.85},0.42) -- ({#1+#3*0.85},0.58);
  \node[echarge, anchor=south] at (#1,0.65) {#2};
}

  \fetframe

  %- shallow trench isolation either side: the source of the stress
  \fill[gray!22] (-2.2,-1.5) rectangle (-0.15,0);
  \fill[gray!22] (8.95,-1.5) rectangle (11.0,0);
  \node[gray!55!black, scale=0.8] at (-1.2,-0.75) {STI};
  \node[gray!55!black, scale=0.8] at (10.0,-0.75) {STI};
  \draw (-2.2,0) -- (11.0,0);

  %- 1. mechanical stress from the trench walls, pushing inwards
  \draw[armygreen, very thick, ->] (-0.1,-2.15) -- (1.3,-2.15);
  \draw[armygreen, very thick, ->] (8.9,-2.15) -- (7.5,-2.15);
  \node[armygreen, scale=0.85] at (4.4,-2.15) {mechanical stress};

  %- 2. the two fields in the channel. E_y starts at the surface rather
  %- than up inside the gate: drawn from y=0.95 it ran straight through
  %- the word GATE and along the gate terminal, which are the two things
  %- in the figure it must not be confused with.
  \draw[blue, very thick, ->] (4.5,-0.05) -- (4.5,-0.75);
  \node[blue, anchor=west, scale=0.85] at (4.6,-0.5) {$E_y$};
  \draw[blue, very thick, ->] (3.3,-1.0) -- (5.5,-1.0);
  \node[blue, anchor=north, scale=0.85] at (4.4,-1.05) {$E_x$};

  %- 3. traps, right at the surface under the gate
  \foreach \x in {3.1,3.5,3.9,4.3,4.7,5.1,5.5} {
    \node[red, scale=0.75] at (\x,0) {$\boldsymbol{\times}$};
  }
  \node[red, anchor=south east, scale=0.85, align=right] at (2.4,1.55)
    {interface and\\ oxide traps};
  %- Both callouts swing round the gate rather than through it. Drawn
  %- straight they crossed the gate outline and landed on the word GATE,
  %- so the figure appeared to be pointing at the electrode when what it
  %- means is the interface underneath it.
  \draw[red, thick, ->] (2.4,1.35) to[out=-90, in=155] (3.05,-0.16);

  %- 4. the drain end, where the lateral field peaks
  \draw[red, very thick, ->] (7.4,1.35) to[out=-90, in=25] (5.75,-0.16);
  \node[red, anchor=south west, scale=0.85, align=left] at (7.5,1.55)
    {hot carriers where\\ the field peaks};

  \node[anchor=north, align=center, scale=0.9] at (4.4,-2.75)
    {Four families of effect, each with its own corner of the device model,\\
     and each taken separately in the sections that follow. Below about\\
     100 nm every one of them is first order.};

\end{tikzpicture}
\end{document}
